\section*{Capítulo \ref{ch:g6:plants-are-producers} --- As plantas são produtoras}

\begin{solution}{exo:g6:plants-are-producers:1}
\omterm{def:g6:plants-are-producers:matter}{Matéria mineral}: a matéria do mundo não vivo --- água, rocha (e
também os gases do ar, os \omterm{prop:g4:what-plants-need:needs}{sais minerais}). \omterm{def:g6:plants-are-producers:matter}{Matéria orgânica}: a
matéria fabricada por \omterm{def:g6:exploring-our-environment:components}{seres vivos} e que compõe o corpo deles ---
madeira, \omterm{def:g3:bones-and-muscles:muscle}{músculo} (e também açúcar, lã, gordura).
\end{solution}

\begin{solution}{exo:g6:plants-are-producers:2}
Ingredientes: água e \omterm{prop:g4:what-plants-need:needs}{sais minerais} do solo, e um gás (o gás
carbônico) do ar. \omterm{prop:g3:balanced-diet:jobs}{Energia}: a \omterm{def:g6:exploring-our-environment:conditions}{luz}, captada pela substância verde.
Produto: \omterm{def:g6:plants-are-producers:matter}{matéria orgânica} --- açúcar e, depois, madeira, \omterm{def:g1:plants-around-us:parts}{folha}, \omterm{def:g1:plants-around-us:parts}{raiz},
\omterm{def:g2:from-seed-to-plant:seed}{semente}. Oficina: a \omterm{def:g1:plants-around-us:parts}{folha} verde, de dia.
\end{solution}

\begin{solution}{exo:g6:plants-are-producers:3}
Mineral: o orvalho, o granito, o sal marinho. Orgânico: a cera de
abelha, o óleo de cozinha --- e a \omterm{def:g1:animals-around-us:covering}{concha} de caracol é fabricação do
caracol, construída com substância mineral que o \omterm{def:g1:animals-around-us:animal}{animal} depositou:
feita por um \omterm{def:g6:exploring-our-environment:components}{ser vivo}, a partir de suprimentos minerais. (Vale
qualquer das duas respostas com esse raciocínio.)
\end{solution}

\begin{solution}{exo:g6:plants-are-producers:4}
Quase todos do ar e da água --- o gás retirado do ar e a água bebida
pelas \omterm{def:g1:plants-around-us:parts}{raízes} ---, montados nas \omterm{def:g1:plants-around-us:parts}{folhas} com a \omterm{prop:g3:balanced-diet:jobs}{energia} da \omterm{def:g6:exploring-our-environment:conditions}{luz}; o solo deu
só pequenas quantidades de \omterm{prop:g4:what-plants-need:needs}{sais minerais}.
\end{solution}

\begin{solution}{exo:g6:plants-are-producers:5}
Ela retira gás carbônico e libera oxigênio. Todo \omterm{def:g6:exploring-our-environment:components}{ser vivo} que respira
expira o primeiro e precisa do segundo --- nós inclusive.
\end{solution}

\begin{solution}{exo:g6:plants-are-producers:6}
O \omterm{ex:g4:plants-without-seeds:bulbtuber}{tubérculo} da batata e o \omterm{ex:g4:plants-without-seeds:bulbtuber}{bulbo} da tulipa (reservas subterrâneas), as
metades gordinhas do feijão (a \omterm{def:g3:animals-through-seasons:active}{reserva de alimento} da \omterm{def:g2:from-seed-to-plant:seed}{semente}), a
\omterm{def:g1:plants-around-us:parts}{raiz} gorda da cenoura, a beterraba cheia de açúcar --- todas reservas
orgânicas depositadas a partir da fabricação do verão.
\end{solution}

\begin{solution}{exo:g6:plants-are-producers:7}
Porque a fabricação funciona com \omterm{def:g6:exploring-our-environment:conditions}{luz}: sem \omterm{def:g6:exploring-our-environment:conditions}{luz} não há \omterm{prop:g3:balanced-diet:jobs}{energia} para
construir \omterm{def:g6:plants-are-producers:matter}{matéria orgânica}, por mais minerais que esperem no solo. A
\omterm{def:g6:exploring-our-environment:conditions}{luz} é o único ingrediente sem substituto --- os sais são tempero, não
combustível.
\end{solution}

\begin{solution}{exo:g6:plants-are-producers:8}
Porque o primeiro elo de uma cadeia precisa colocar \omterm{def:g6:plants-are-producers:matter}{matéria orgânica}
nova no mundo vivo, e os \omterm{def:g1:animals-around-us:animal}{animais} são \omterm{def:g5:food-webs:roles}{consumidores}: eles só conseguem
transformar \omterm{def:g6:plants-are-producers:matter}{matéria orgânica} já fabricada por outra pessoa. Só as
produtoras fabricam essa matéria a partir de suprimentos minerais ---
então o primeiro elo é sempre verde.
\end{solution}

\begin{solution}{exo:g6:plants-are-producers:9}
Pão: trigo --- as \omterm{prop:g3:flowers-fruits-seeds:transform}{sementes} do pé de trigo; a trilha termina na \omterm{def:g1:plants-around-us:parts}{folha}
do trigo. Manteiga: creme --- \omterm{def:g2:animals-grow-up:born}{leite} --- a vaca --- o capim; a trilha
termina na \omterm{def:g1:plants-around-us:parts}{folha} do capim. Queijo: \omterm{def:g2:animals-grow-up:born}{leite} --- a vaca --- o capim de
novo. Três alimentos, três (ou duas) \omterm{def:g1:plants-around-us:parts}{folhas}, todas na \omterm{def:g6:exploring-our-environment:conditions}{luz}.
\end{solution}

\begin{solution}{exo:g6:plants-are-producers:10}
O feno é o alimento da vaca --- \omterm{def:g6:plants-are-producers:matter}{matéria orgânica} que ela transforma no
corpo dela. O composto não é o alimento da \omterm{def:g1:plants-around-us:plant}{planta}: ele reabastece os
sais \emph{minerais} do solo enquanto apodrece. A \omterm{def:g1:plants-around-us:plant}{planta} então fabrica
a própria \omterm{def:g6:plants-are-producers:matter}{matéria orgânica} a partir desses minerais, do ar, da água e
da \omterm{def:g6:exploring-our-environment:conditions}{luz}. A vaca come refeições; a \omterm{def:g1:plants-around-us:plant}{planta} fabrica as dela.
\end{solution}

\begin{solution}{exo:g6:plants-are-producers:11}
Iluminado, a fabricação diurna das \omterm{def:g1:plants-around-us:plant}{plantas} libera oxigênio e consome
gás carbônico, equilibrando a \omterm{def:g4:breathing:movements}{respiração} dos \omterm{prop:g3:sorting-living-things:groups}{peixes} --- a troca de
gases funciona nos dois sentidos e o vidro fecha o ciclo. No escuro, a
fabricação para enquanto todos os moradores (as \omterm{def:g1:plants-around-us:plant}{plantas} inclusive)
continuam respirando: o oxigênio acaba, o gás carbônico se acumula e o
mundo fechado sufoca.
\end{solution}

\begin{solution}{exo:g6:plants-are-producers:12}
Certa: a matéria da árvore não veio da terra --- a pesagem dele provou
isso para sempre. Incompleta: a água era só um dos ingredientes de
volume; ele não tinha como saber do gás retirado do ar, que fornece
boa parte da substância da madeira. O acréscimo moderno: o gás
carbônico do ar (mais os pequenos \omterm{prop:g4:what-plants-need:needs}{sais minerais} do solo) --- água e ar
juntos, montados pela \omterm{def:g6:exploring-our-environment:conditions}{luz}.
\end{solution}

\begin{solution}{pb:g6:plants-are-producers:1}
\textbf{1.} As toneladas são \omterm{def:g6:plants-are-producers:matter}{matéria orgânica} fabricada pelas próprias
\omterm{def:g1:plants-around-us:plant}{plantas} a partir de suprimentos minerais --- água, o gás carbônico do
ar, os sais do solo --- com a \omterm{def:g6:exploring-our-environment:conditions}{luz} como \omterm{prop:g3:balanced-diet:jobs}{energia}. Os sacos reabastecem
os minerais; o ar e a chuva entregam o volume; as \omterm{def:g1:plants-around-us:parts}{folhas} fazem a
fabricação.

\textbf{2.} Batata: órgão de reserva subterrâneo (\omterm{ex:g4:plants-without-seeds:bulbtuber}{tubérculo}).
Cenoura: \omterm{def:g1:plants-around-us:parts}{raiz} de reserva. \omterm{def:g2:from-seed-to-plant:seed}{Semente} de feijão: a \omterm{def:g3:animals-through-seasons:active}{reserva de alimento} da
\omterm{def:g2:from-seed-to-plant:seed}{semente}, embalada para a \omterm{def:g1:plants-around-us:plant}{planta} seguinte. \omterm{def:g1:plants-around-us:parts}{Folha} de alface: o próprio
órgão em atividade --- a oficina que capta \omterm{def:g6:exploring-our-environment:conditions}{luz}, comida fresca.

\textbf{3.} O fornecedor de verdade é a \omterm{def:g6:exploring-our-environment:conditions}{luz}: cada grama da produção
foi construída com a \omterm{prop:g3:balanced-diet:jobs}{energia} da \omterm{def:g6:exploring-our-environment:conditions}{luz}, de dia, nas \omterm{def:g1:plants-around-us:parts}{folhas}.

\textbf{4.} Pelos \omterm{prop:g4:what-plants-need:needs}{sais minerais}: as colheitas levam esses sais
embora e, sem reabastecimento, as \omterm{def:g1:plants-around-us:plant}{plantas} ficariam com falta de
minerais como a hortelã na areia --- pequenos em volume,
indispensáveis em função.

\textbf{5.} A experiência do armário escuro (o prisioneiro pálido): o
pano removeu a \omterm{def:g6:exploring-our-environment:conditions}{luz}, e uma semana já aparece --- pálidas, esticadas,
com fabricação fraca.

\textbf{6.} Os \omterm{ex:g4:plants-without-seeds:bulbtuber}{tubérculos} são enchidos pelas ramas: as \omterm{def:g1:plants-around-us:parts}{folhas}
fabricam açúcar de dia, e a sobra é canalizada para baixo e
depositada nos \omterm{ex:g4:plants-without-seeds:bulbtuber}{tubérculos} que incham. \omterm{def:g1:plants-around-us:parts}{Folhas} frondosas em atividade
querem dizer uma sobra grande --- boa rama, bom \omterm{ex:g4:plants-without-seeds:bulbtuber}{tubérculo}.

\textbf{7.} A água. A fabricação precisa de todos os ingredientes ao
mesmo tempo: a \omterm{def:g6:exploring-our-environment:conditions}{luz} fornece \omterm{prop:g3:balanced-diet:jobs}{energia}, não substância --- com as \omterm{def:g1:plants-around-us:parts}{raízes}
secas, não há com o que construir, por mais brilhante que esteja o céu.

\textbf{8.} As \omterm{def:g1:plants-around-us:parts}{raízes} enterradas são \omterm{def:g6:plants-are-producers:matter}{matéria orgânica} devolvida ao
solo: apodrecendo, elas vão liberar os minerais delas para as
lavouras seguintes --- a regra da devolução da
\cref{prop:g5:people-and-nature:lasting}, paga em matéria vegetal. O
desperdício seria queimá-las.

\textbf{9.} Pé de batata $\rightarrow$ porco $\rightarrow$ a família.
Toda a \omterm{def:g6:plants-are-producers:matter}{matéria orgânica} dessa cadeia foi fabricada primeiro nas
\omterm{def:g1:plants-around-us:parts}{folhas} do pé de batata --- o porco e a família só a transformaram.

\textbf{10.} Ingredientes: água e \omterm{prop:g4:what-plants-need:needs}{sais minerais} do solo, gás
carbônico do ar --- montados em matéria de \omterm{def:g1:plants-around-us:parts}{folha}. O papel exato do
sol: \omterm{prop:g3:balanced-diet:jobs}{energia}, e não substância. A salada dela é água, ar e terra
\emph{construídos por} \omterm{def:g6:exploring-our-environment:conditions}{luz} do sol --- o crédito poético é merecido, a
lista de ingredientes é que está errada.

\textbf{11.} Os \omterm{def:g5:food-webs:roles}{decompositores} --- a vida do solo, das minhocas e dos
tatuzinhos até \omterm{def:g6:exploring-our-environment:components}{seres vivos} pequenos demais para enxergar --- se
alimentam das \omterm{def:g1:plants-around-us:parts}{folhas} mortas e quebram a \omterm{def:g6:plants-are-producers:matter}{matéria orgânica} delas rumo à
\omterm{def:g6:plants-are-producers:matter}{matéria mineral} de novo, reabastecendo o solo de que as próximas
lavouras vão tirar.

\textbf{12.} Por exemplo: ``As produtoras da horta constroem \omterm{def:g6:plants-are-producers:matter}{matéria
orgânica} --- a colheita --- a partir de \omterm{def:g6:plants-are-producers:matter}{matéria mineral} do solo, da
água e do ar, usando a \omterm{def:g6:exploring-our-environment:conditions}{luz} como \omterm{prop:g3:balanced-diet:jobs}{energia} delas.''
\end{solution}
